\documentclass[12 pt, letterpaper]{hmcpset}
\usepackage{hw-preamble}

\name{Jayden Thadani}
\class{MATH 82 --- Differential Equations}
\assignment{Homework 3}
\duedate{}

\begin{document}

\begin{problem}[1:]
	Consider the initial value problem
	\[
		y' = 1 - t + y, \quad y(0) = 1.
	\]
	\begin{enumerate}
		\item Use Euler's method with step-size $h = 0.2$ to estimate the value of $y(1)$. Then, repeat for $h = 0.1$.
		\item Find the solution to the differential equation analytically and use it to determine the exact value of $y(1)$.
		\item Verify that halving the step-size $h$ roughly halves the error in your numerical approximations of $y(1)$.
	\end{enumerate}
\end{problem}

\begin{solution}
	$y_{n}$ as defined below by Euler's method approximates $y$.
	\[
		y_{n} = y_{n - 1}' h + y_{n - 1}.
	\]
	Specifically, $y_{j} \approx y(jh)$ and $y_{j}' \approx y'(jh)$. We implement this recurrence in the following Python program to solve differential equations using Euler's method.
\begin{verbatim}
def EulersMethod(y_0, t_0, t_f, h, ydot):
    y = [y_0]
    t = [t_0]
    ydots = []
    while (t[-1] < t_f):
        ydots.append(ydot(y[-1], t[-1]))
        y.append(ydots[-1] * h + y[-1])
        t.append(t[-1] + h)
    return y, ydots	
\end{verbatim}
	With a function \texttt{ydot} that takes input \texttt{y} and \texttt{t}, returning \texttt{1 - t + y}, we can numerically solve the differential equation.

	\begin{enumerate}
		\item
			\[
				\boxed{
					y \approx \begin{cases}
					3.48832 & \text{for } h = 0.2 \\
					3.59374 & \text{for } h = 0.1
				\end{cases}
			}
			\]

			Using the Python code --- executing \texttt{EulersMethod(1, 0, 1, 0.2, ydot)} --- we get the following lists for $h = 0.2$.
			\begin{align*}
				y_{0} & = 1 & y_{0}' & = 2 \\
				y_{1} & = 1.4 & y_{1}' & = 2.2 \\
				y_{2} & = 1.84 & y_{2}' & = 2.44 \\
				y_{3} & = 2.328 & y_{3}' & = 2.728 \\
				y_{4} & = 2.8736 & y_{4} & = 3.0736 \\
				y_{5} & = 3.48832
			\end{align*}
			Similarly, for $h = 0.1$ we execute \texttt{EulersMethod(1, 0, 1, 0,1, ydot)} and get $y_{10} \approx 3.59374$ (rounded to five decimal places).
		\item
			\[
				\boxed{
					y(1) = 1 + e.
				}
			\]
			The equation is just a first order linear differential equation. We can solve it by just noticing that the integrating factor has to be $e^{-t}$. This gives us
			\[
				(y e^{-t})' = e^{-t} (1 - t).
			\]
			Note that the right hand side is just the $t$-derivative of $e^{-t} t$ and thus, we have
			 \[
				 (y e^{-t})' = (e^{-t} t)'.
			\]
			Then by the fundamental theorem of calculus,
			\[
				y e^{-t} = e^{-t} t + C.
			\]
			$C$ can be determined by considering the initial condition $y(0) = 1$ which forces us to have
			\[
				1 e^{0} = 0 e^{0} + C
			\]
			which is just $C = 1$. Thus,
			\[
				y(t) = t + e^{t}.
			\]
			and $y(1) = 1 + e$.
		\item
			The error is about $-0.22996$ for $h = 0.2$ and about half the magnitude --- it's approximately $0.12088$ for $h = 0$.
	\end{enumerate}
\end{solution}

\pagebreak

\begin{problem}[2:]
	Consider the differential equation $y' = \cos{t} - y$.
	\begin{enumerate}
		\item Find the general solution to this ordinary differential equation.
		\item Generate a direction field and some solution curves for this differential equation with different initial conditions.
		\item Based on your graphical exploration in part (b), conjecture what happens to $y(t)$ as $t$ tends to infinity. Corroborate that conjecture using the general solution you found in part (a).
	\end{enumerate}
\end{problem}

\pagebreak

\begin{problem}[3:]
	\begin{enumerate}
		\item Sketch a direction field for the differential equation
			\[
				y' = - \frac{x}{y}
			\]
			for $x \in [-3, 3]$ and $y \in [-3, 3]$.
		\item If you were given the sketch of the direction field from part (a), but were not given the equation, how would you know that the differential equation was non-autonomous?
	\end{enumerate}
\end{problem}

\pagebreak

\begin{problem}[4:]
	You have been hired by a fishery to do some preliminary work on modelling a fish population under various harvesting strategies. They have asked you to start with this mathematical model that accounts for logistic growth, death, and harvesting:
	\[
		\frac{dP}{dt} = r P(1 - P/K) - \alpha P - H(P).
	\]
	Here, $P(t)$ represents the fishery biomass (total mass of fish). The constant $r$ relates to the growth rate of your population, and $K$ is the carrying capacity of your fishery --- the maximum biomass that the fishery can sustain. The term $-\alpha P$ accounts for the natural death of your fish. Assume $r$, $K$, $\alpha$ and $P(t)$ are all positive.

	The function $H(t, P)$ describes the harvesting of the fish. The fishery wants to understand what would happen to the population of fish under different scenarios.

	\begin{enumerate}
		\item Suppose we use kilograms as the units for $P$ and days as the units for $t$. What must be the units for $r$, $\alpha$, and $K$?
		\item Notice that the problem assumes the harvesting rate of the fish is time invariant since $H(P)$ isn't directly a function of $t$. Here are two natural choices for $H$ in this case: the fishery could harvest at a constant rate or it could harvest at a rate that is proportional to the fishery biomass. In each of these cases, calculate the long-term population of the fish, and the long-term harvesting rate.
		\item In the case where $H = H_{0}$ is constant, solve for $P(t)$ analytically. In addition, numerically approximate $P(t)$ and verify that your analytic answer agrees with your numerical one.
	\end{enumerate}
\end{problem}

\end{document}
